\appendix{Qualitative Results (ImProver)}
\label{app:qual_improver}


In this section, we provide additional qualitative examples demonstrating the improvements \methodname{} achieves in proof optimization.

\vspace{15px}
\paragraph{Compfiles: Length Optimization} 

Consider \autoref{fig:ex1-comp-len},  a lemma from the 2022 IMO Question 2 (Compfiles) that we optimize for length. \methodname{} halves thr proof from 12 tactics to 6. Here, \methodname{} makes multiple nontrivial optimizations, such as eliminating the \texttt{h2'} and \texttt{h4} and \texttt{hxw} hypotheses, as well as fully generating proof terms for specific rewrites and other tactics.
\vspace{15px}
\paragraph{Compfiles: Declarativity Optimization} 



Consider \autoref{fig:compfiles-read}, in which a lemma from the 2019 IMO problem 1 (from the Compfiles dataset) is optimized for declarativity. This introduces multiple new hypotheses, which generalize a \texttt{linear\_property} of the functions, and then reuses and instantiates that (and others, too) hypothesis throughout the proof, creating a significantly more declarative proof.


\setlength{\columnsep}{0.2cm}
\begin{figure}[ht]
\begin{multicols}{2}

{
\textbf{Original (human-written)}
\begin{scriptsize}
\begin{lstlisting}[basicstyle=\scriptsize\ttfamily]
lemma additive_to_int_linear (f : ℤ → ℤ) (h: ∀ (x y : ℤ), f (x + y) = f x + f y):
   ∃ c, ∀ a, f a = c * a  := by
  let g := AddMonoidHom.toIntLinearMap <| AddMonoidHom.mk' f h
  refine ⟨f 1, fun a => ?_⟩
  change g a = g 1 * a
  rw [mul_comm, ← smul_eq_mul, ← LinearMap.map_smul, smul_eq_mul, mul_one]
\end{lstlisting}
\end{scriptsize}
}
\vfill\null
\columnbreak
{
\textbf{\methodname{} (declarativity-optimized)}
\begin{scriptsize}
\begin{lstlisting}[basicstyle=\scriptsize\ttfamily]
lemma additive_to_int_linear (f : ℤ → ℤ) (h: ∀ (x y : ℤ), f (x + y) = f x + f y):
   ∃ c, ∀ a, f a = c * a   := by
  let g := AddMonoidHom.toIntLinearMap <| AddMonoidHom.mk' f h
  have linear_property : ∀ a, f a = g a := by
    intro a
    rfl
  have g_smul : ∀ a, g a = g 1 * a := by
    intro a
    rw [mul_comm, ← smul_eq_mul, ← LinearMap.map_smul, smul_eq_mul, mul_one]
  refine ⟨f 1, fun a => ?_⟩
  have f_eq_g : f a = g a := linear_property a
  have g_a_eq : g a = g 1 * a := g_smul a
  rw [f_eq_g, linear_property 1, g_a_eq]
\end{lstlisting}
\end{scriptsize}
}
\end{multicols}
\vspace{-2.5em}
\caption{Optimizing a lemma from IMO 2019 P1 for declarativity}
\label{fig:compfiles-read}
\end{figure}

\vspace{15px}
\paragraph{Compfiles: Mixed Optimization}

Consider \autoref{fig:compfiles-mix}, in which a lemma from the 2023 USAMO problem 2 (from the Compfiles dataset) is optimized for mixed declarativity and length. This introduces a new hypothesis, which declares a powerful intermediate lemma, which is then applied to solve the problem. Moreover, this declarativity is introduced in such a way that it makes the proof more concise than the original, with $3$ tactics rather than $5$.


\setlength{\columnsep}{0.2cm}
\begin{figure}[ht]
\begin{multicols}{2}

{
\textbf{Original (human-written)}
\begin{scriptsize}
\begin{lstlisting}[basicstyle=\scriptsize\ttfamily]
lemma lemma_3 {a b c : ℝ+} (h : a = b + c) : c < a  := by
  rw [h]
  obtain ⟨b, hb⟩ := b
  obtain ⟨c, hc⟩ := c
  rw [←Subtype.coe_lt_coe, Positive.coe_add]
  exact lt_add_of_pos_left c hb
\end{lstlisting}
\end{scriptsize}
}
\vfill\null
\columnbreak
{
\textbf{\methodname{} (mix-optimized)}
\begin{scriptsize}
\begin{lstlisting}[basicstyle=\scriptsize\ttfamily]
lemma lemma_3 {a b c : ℝ+} (h : a = b + c) : c < a    := by
  have : ↑c < ↑b + ↑c := lt_add_of_pos_left c.val b.property
  rw [h, ←Subtype.coe_lt_coe, Positive.coe_add]
  exact this
\end{lstlisting}
\end{scriptsize}
}
\end{multicols}
\vspace{-2.5em}
\caption{Optimizing a lemma from USAMO 2023 P2 for mixed declarativity/length}
\label{fig:compfiles-mix}
\end{figure}

\vspace{15px}
\paragraph{MIL: Length Optimization} 

Consider \autoref{fig:mil-len}, which optimizes an exercise solution from MIL Chapter 8, Section 1 (Group theory) for length, modifying the proof structure and introducing proof terms into the structure of the proof to shorten it from $9$ tactic invocations to $7$.

\setlength{\columnsep}{0.2cm}
\begin{figure}[ht]
\begin{multicols}{2}

{
\textbf{Original (human-written)}
\begin{scriptsize}
\begin{lstlisting}[basicstyle=\scriptsize\ttfamily]
example (φ : G →* H) (ψ : H →* K) (S : Subgroup G) :
    map (ψ.comp φ) S = map ψ (S.map φ)  := by
  ext x
  simp only [mem_map]
  constructor
  · rintro ⟨y, y_in, hy⟩
    exact ⟨φ y, ⟨y, y_in, rfl⟩, hy⟩
  · rintro ⟨y, ⟨z, z_in, hz⟩, hy⟩
    use z, z_in
    calc ψ.comp φ z = ψ (φ z) := rfl
    _               = ψ y := by congr
\end{lstlisting}
\end{scriptsize}
}
\vfill\null
\columnbreak
{
\textbf{\methodname{} (length-optimized)}
\begin{scriptsize}
\begin{lstlisting}[basicstyle=\scriptsize\ttfamily]
example (φ : G →* H) (ψ : H →* K) (S : Subgroup G) :
    map (ψ.comp φ) S = map ψ (S.map φ)   := by
  ext x
  simp only [mem_map]
  constructor
  rintro ⟨y, y_in, hy⟩; exact ⟨φ y, ⟨y, y_in, rfl⟩, hy⟩
  rintro ⟨y, ⟨z, z_in, hz⟩, hy⟩; exact ⟨z, z_in, (congr_arg ψ hz).trans hy⟩
\end{lstlisting}
\end{scriptsize}
}
\end{multicols}
\vspace{-2.5em}
\caption{Optimizing a lemma from the solutions of MIL CH08 S01 for length}
\label{fig:mil-len}
\end{figure}

\vspace{15px}
\paragraph{MIL Declarativity Optimization.}
\label{ex2}
Consider \autoref{fig:ex2-MIL-read}, a result from MIL that we optimize for the declarative metric. This original proof carried a score of $0$, as it has no \texttt{have} statements. In comparison, after applying \methodname{}, the model explicitly defines hypotheses for use in the latter half of the proof; these hypotheses can easily be converted into standalone lemmas for reuse.



\setlength{\columnsep}{0.2cm}
\begin{figure}[ht]
\begin{multicols}{2}

{
\textbf{Original (human-written)}
\begin{scriptsize}
\begin{lstlisting}[basicstyle=\scriptsize\ttfamily]
def iso₁ [Fintype G] (h : Disjoint H K) (h' : card G = card H * card K) 
: K ≃* G / H  := by
  apply MulEquiv.ofBijective ((QuotientGroup.mk' H).restrict K)
  rw [bijective_iff_injective_and_card]
  constructor
  · rw [← ker_eq_bot_iff, (QuotientGroup.mk' H).ker_restrict K]
    simp [h]
  · symm
    exact aux_card_eq h'
\end{lstlisting}
\end{scriptsize}
}
\vfill\null
\columnbreak
{
\textbf{\methodname{} (declarativity-optimized)}
\begin{scriptsize}
\begin{lstlisting}[basicstyle=\scriptsize\ttfamily]
def iso₁ [Fintype G] (h : Disjoint H K) (h' : card G = card H * card K) 
: K ≃* G / H   := by
  have injectivity : Function.Injective ((QuotientGroup.mk' H).restrict K) := by
    rw [← ker_eq_bot_iff, (QuotientGroup.mk' H).ker_restrict K]
    simp [h]
  have card_eq : card (G / H) = card K := by
    exact aux_card_eq h'
  apply MulEquiv.ofBijective ((QuotientGroup.mk' H).restrict K)
  rw [bijective_iff_injective_and_card]
  constructor
  exact injectivity
  symm
  exact card_eq
\end{lstlisting}
\end{scriptsize}
}
\end{multicols}
\vspace{-2.5em}
\caption{Optimizing a group-theoretic result from MIL Chapter 8 Section 1 for declarativity.}
\label{fig:ex2-MIL-read}
\vspace{-1em}
\end{figure}



\vspace{15px}
\paragraph{MIL: Mixed Optimization}

Consider \autoref{fig:mil-mix}, which optimizes an exercise solution from MIL Chapter 8, Section 1 (Group theory) for mixed length/declarativity. \methodname{} significantly modifies the structure of the proof, applying more general-purpose tactics like \texttt{simpa} to simplify the proof (via eliminating the need for manual step-by-step calculations). Moreover, the proof introduces declarativity by a intermediate result, which is used by the \texttt{simpa} call to finish the proof in an efficient and declarative manner.


\setlength{\columnsep}{0.2cm}
\begin{figure}[ht]
\begin{multicols}{2}

{
\textbf{Original (human-written)}
\begin{scriptsize}
\begin{lstlisting}[basicstyle=\scriptsize\ttfamily]
lemma eq_bot_iff_card {G : Type*} [Group G] {H : Subgroup G} [Fintype H] :
    H = ⊥ ↔ card H = 1  := by
  suffices (∀ x ∈ H, x = 1) ↔ ∃ x ∈ H, ∀ a ∈ H, a = x by
    simpa [eq_bot_iff_forall, card_eq_one_iff]
  constructor
  · intro h
    use 1, H.one_mem
  · rintro ⟨y, -, hy'⟩ x hx
    calc x = y := hy' x hx
    _      = 1 := (hy' 1 H.one_mem).symm
\end{lstlisting}
\end{scriptsize}
}
\vfill\null
\columnbreak
{
\textbf{\methodname{} (mix-optimized)}
\begin{scriptsize}
\begin{lstlisting}[basicstyle=\scriptsize\ttfamily]
lemma eq_bot_iff_card {G : Type*} [Group G] {H : Subgroup G} [Fintype H] :
    H = ⊥ ↔ card H = 1     := by
  have : (∀ x ∈ H, x = 1) ↔ ∃ x ∈ H, ∀ a ∈ H, a = x :=
    ⟨λ h => ⟨1, H.one_mem, h⟩, λ ⟨y, _, hy'⟩ x hx => (hy' 1 H.one_mem).symm ▸ hy' x hx⟩
  simpa [eq_bot_iff_forall, card_eq_one_iff] using this
\end{lstlisting}
\end{scriptsize}
}
\end{multicols}
\vspace{-2.5em}
\caption{Optimizing a lemma from MIL CH08 S01 solution for mixed declarativity/length}
\label{fig:mil-mix}
\end{figure}


\vspace{15px}
\paragraph{Mathlib: Length Optimization} 

Consider \autoref{fig:mathlib-len}, which optimizes a theorem in algebraic topology from mathlib for length, eliminating \texttt{simp} calls and combining tactics to shorten it from $3$ tactic invocations to $1$.


\setlength{\columnsep}{0.2cm}
\begin{figure}[ht]
\begin{multicols}{2}

{
\textbf{Original (human-written)}
\begin{scriptsize}
\begin{lstlisting}[basicstyle=\scriptsize\ttfamily,literate={π}{{\ensuremath{\mathrm{\mathnormal{\pi}}}}}1 {ₘ}{{\ensuremath{_m}}}1 {←}{{\ensuremath{\leftarrow}}}1 {lbrak}{{$\llbracket$}}1 {rbrak}{{$\rrbracket$}}1]
/-- If `f(p(t) = g(q(t))` for two paths `p` and `q`, then the induced path homotopy classes
`f(p)` and `g(p)` are the same as well, despite having a priori different types -/
theorem heq_path_of_eq_image : HEq ((πₘ f).map lbrak p rbrak) ((πₘ g).map lbrak q rbrak)  := by
  simp only [map_eq, ← Path.Homotopic.map_lift]; apply Path.Homotopic.hpath_hext; exact hfg
\end{lstlisting}
\end{scriptsize}
}
\vfill\null
\columnbreak
{
\textbf{\methodname{} (length-optimized)}
\begin{scriptsize}
\begin{lstlisting}[basicstyle=\scriptsize\ttfamily,literate={π}{{\ensuremath{\mathrm{\mathnormal{\pi}}}}}1 {ₘ}{{\ensuremath{_m}}}1 {lbrak}{{$\llbracket$}}1 {rbrak}{{$\rrbracket$}}1]
/-- If `f(p(t) = g(q(t))` for two paths `p` and `q`, then the induced path homotopy classes
`f(p)` and `g(p)` are the same as well, despite having a priori different types -/
theorem heq_path_of_eq_image : HEq ((πₘ f).map lbrak p rbrak) ((πₘ g).map lbrak q rbrak)    := by
  exact Path.Homotopic.hpath_hext hfg
\end{lstlisting}
\end{scriptsize}
}
\end{multicols}
\vspace{-2.5em}
\caption{Optimizing a theorem from \texttt{Mathlib/FundamentalGroupoid/InducedMaps} for length}
\vspace{-1em}
\label{fig:mathlib-len}
\end{figure}


\vspace{15px}
\paragraph{Mathlib: Declarativity Optimization} 

Consider \autoref{fig:mathlib-read}, a theorem from Mathlib that we optimize for declarativity. 

This original proof carried a score of $0$, as it does not contain any declarative statements. It is concise and efficient, however, it is difficult to understand and read. 

After optimizing for declarativity, we see that the model did not change the structure of the proof. Rather, it added an intermediate declaration so that users can better understand the state after the \texttt{convert}. This intermediate tactic greatly helps in the understandability and clarity of the proof.


\setlength{\columnsep}{0.2cm}
\begin{figure}[ht]
\begin{multicols}{2}

{
\textbf{Original (human-written)}
\begin{scriptsize}
\begin{lstlisting}[basicstyle=\scriptsize\ttfamily]
/-- Another version of `simply_connected_iff_paths_homotopic` -/
theorem simply_connected_iff_paths_homotopic' {Y : Type*} [TopologicalSpace Y] :
    SimplyConnectedSpace Y ↔
      PathConnectedSpace Y ∧ ∀ {x y : Y} (p₁ p₂ : Path x y), Path.Homotopic p₁ p₂  := by
  convert simply_connected_iff_paths_homotopic (Y := Y)
  simp [Path.Homotopic.Quotient, Setoid.eq_top_iff]; rfl
\end{lstlisting}
\end{scriptsize}
}
\vfill\null
\columnbreak
{
\textbf{\methodname{} (declarativity-optimized)}
\begin{scriptsize}
\begin{lstlisting}[basicstyle=\scriptsize\ttfamily]
/-- Another version of `simply_connected_iff_paths_homotopic` -/
theorem simply_connected_iff_paths_homotopic' {Y : Type*} [TopologicalSpace Y] :
    SimplyConnectedSpace Y ↔
      PathConnectedSpace Y ∧ ∀ {x y : Y} (p₁ p₂ : Path x y), Path.Homotopic p₁ p₂   := by
  convert simply_connected_iff_paths_homotopic (Y := Y)
  have h1 : ∀ x y : Y, Subsingleton (Path.Homotopic.Quotient x y) ↔ ∀ {p₁ p₂ : Path x y}, Path.Homotopic p₁ p₂ := by
    intro x y
    simp [Path.Homotopic.Quotient, Setoid.eq_top_iff]
    rfl
  simp only [h1]
\end{lstlisting}
\end{scriptsize}
}
\end{multicols}
\vspace{-2.5em}
\caption{Optimizing a theorem from \texttt{Mathlib/FundamentalGroupoid/SimplyConnected} for declarativity}
\vspace{-1em}
\label{fig:mathlib-read}
\end{figure}




\vspace{15px}
\paragraph{Mathlib: Mixed Optimization}


Consider \autoref{fig:mathlib-mix}, a theorem from Mathlib that we optimize for mixed length/declarativity. 

We observe that \methodname{} applies more complex tactics such as \texttt{all\_goals} and \texttt{split\_ifs} to significantly decrease the number of tactics in the proofs, while maintaining the overall structure.


\setlength{\columnsep}{0.2cm}
\begin{figure}[ht]
\begin{multicols}{2}

{
\textbf{Original (human-written)}
\begin{scriptsize}
\begin{lstlisting}[basicstyle=\scriptsize\ttfamily]
theorem trans_refl_reparam (p : Path x₀ x₁) :
    p.trans (Path.refl x₁) =
      p.reparam (fun t => ⟨transReflReparamAux t, transReflReparamAux_mem_I t⟩) (by continuity)
        (Subtype.ext transReflReparamAux_zero) (Subtype.ext transReflReparamAux_one)  := by
      p.reparam (fun t => ⟨transReflReparamAux t, transReflReparamAux_mem_I t⟩) (by continuity)
        (Subtype.ext transReflReparamAux_zero) (Subtype.ext transReflReparamAux_one) := by
  ext
  unfold transReflReparamAux
  simp only [Path.trans_apply, not_le, coe_reparam, Function.comp_apply, one_div, Path.refl_apply]
  split_ifs
  · rfl
  · rfl
  · simp
  · simp
\end{lstlisting}
\end{scriptsize}
}
\vfill\null
\columnbreak
{
\textbf{\methodname{} (mix-optimized)}
\begin{scriptsize}
\begin{lstlisting}[basicstyle=\scriptsize\ttfamily]
theorem trans_refl_reparam (p : Path x₀ x₁) :
    p.trans (Path.refl x₁) =
      p.reparam (fun t => ⟨transReflReparamAux t, transReflReparamAux_mem_I t⟩) (by continuity)
        (Subtype.ext transReflReparamAux_zero) (Subtype.ext transReflReparamAux_one)   := by
      p.reparam (fun t => ⟨transReflReparamAux t, transReflReparamAux_mem_I t⟩) (by continuity)
        (Subtype.ext transReflReparamAux_zero) (Subtype.ext transReflReparamAux_one)  := by
  ext t
  simp only [Path.trans_apply, Path.refl_apply, coe_reparam, Function.comp_apply]
  unfold transReflReparamAux
  split_ifs
  all_goals { simp }
\end{lstlisting}
\end{scriptsize}
}
\end{multicols}
\vspace{-2.5em}
\caption{Optimizing a theorem from \texttt{Mathlib/FundamentalGroupoid/SimplyConnected} for mixed length/declarativity}
\vspace{-1em}
\label{fig:mathlib-mix}
\end{figure}



\vspace{15px}
\paragraph{Full Proof Generation (MIL).}
\label{ex3}
We analyze the application of \methodname{} to neural theorem proving in the MIL example from \autoref{fig:ex3-MIL-completion}. This theorem relating to group theory originally has no proof, however, \methodname{} generates one from scratch. This generated proof is verified to be correct by Lean, utilizing all the included hypotheses as well as a retrieved mathlib theorem.




\setlength{\columnsep}{0.2cm}
\begin{figure}[ht]
\begin{multicols}{2}

{
\textbf{Original (human-written)}
\begin{scriptsize}
\begin{lstlisting}[basicstyle=\scriptsize\ttfamily]
example (φ : G →* H) (S T : Subgroup H) (hST : S ≤ T) : comap φ S ≤ comap φ T  := by
  sorry
\end{lstlisting}
\end{scriptsize}
}
\vfill\null
\columnbreak
{
\textbf{\methodname{} (completeness-optimized)}
\begin{scriptsize}
\begin{lstlisting}[basicstyle=\scriptsize\ttfamily]
example (φ : G →* H) (S T : Subgroup H) (hST : S ≤ T) : comap φ S ≤ comap φ T   := by
  intro g
  simp only [mem_comap]
  intro hS
  exact hST hS
\end{lstlisting}
\end{scriptsize}
}
\end{multicols}
\vspace{-2.5em}
\caption{Solving a group theorem exercise from MIL Chapter 8 Section 1.}
\vspace{-1em}
\label{fig:ex3-MIL-completion}
\end{figure}


\vspace{15px}
\paragraph{Full Proof Generation (MiniF2F).}

We analyze the application of \methodname{} to neural theorem proving in the MiniF2F example from \autoref{fig:ex3-F2F-completion}. This theorem relating to number theory originally has no proof, however, \methodname{} generates one from scratch by calculating step by step, applying theorems retrieved from Mathlib, as well as general tactics like \texttt{simp}.



\setlength{\columnsep}{0.2cm}
\begin{figure}[ht]
\begin{multicols}{2}

{
\textbf{Original (human-written)}
\begin{scriptsize}
\begin{lstlisting}[basicstyle=\scriptsize\ttfamily]
theorem mathd_numbertheory_345 :
  (2000 + 2001 + 2002 + 2003 + 2004 + 2005 + 2006) % 7 = 0 := by sorry

\end{lstlisting}
\end{scriptsize}
}
\vfill\null
\columnbreak
{
\textbf{\methodname{} (completeness-optimized)}
\begin{scriptsize}
\begin{lstlisting}[basicstyle=\scriptsize\ttfamily]

theorem mathd_numbertheory_345 :
  (2000 + 2001 + 2002 + 2003 + 2004 + 2005 + 2006) % 7 = 0   := by
  calc
    (2000 + 2001 + 2002 + 2003 + 2004 + 2005 + 2006) % 7 = (0 + 1 + 2 + 3 + 4 + 5 + 6) % 7 := by
      simp [Nat.add_mod, Nat.mod_eq_zero_of_dvd, Nat.dvd_refl, Nat.mod_add_mod, Nat.mod_self]
    _ = 21 % 7 := by simp
    _ = 0 := by simp
\end{lstlisting}
\end{scriptsize}
}
\end{multicols}
\vspace{-2.5em}
\caption{Solving a number theoretic theorem from MiniF2F.}
\vspace{-1em}
\label{fig:ex3-F2F-completion}
\end{figure}



